\documentclass[journal]{IEEEtran}

\usepackage[english]{babel}
\usepackage{amsmath}
\usepackage{graphicx}
\usepackage{siunitx}
\usepackage{booktabs}
\usepackage{orcidlink}
\usepackage{xcolor} 
\usepackage{cite}
\usepackage{hyperref} % keep last


\sisetup{table-format=1.4}

\title{\textcolor{red}{Quantitative Benchmarking of Multispectral Sentinel-2 Super-Resolution with SEN2NEON}}

\markboth{IEEE GEOSCIENCE AND REMOTE SENSING LETTERS}{Donike \MakeLowercase{\textit{et al.}}}



\author{Simon~Donike\orcidlink{0000-0002-4440-3835},
        Cesar~Aybar\orcidlink{0000-0003-2745-9535},
        Julio~Contreras\orcidlink{0009-0001-5408-7055},
        and~Luis~Gómez\textendash Chova\orcidlink{0000-0003-3924-1269},~\IEEEmembership{Senior Member,~IEEE},%
\thanks{All authors are with the Image Processing Laboratory (IPL), University of Valencia, Spain. Corresponding author: simon.donike@uv.es.}}


\begin{document}
\maketitle

\begin{abstract}
The super-resolution (SR) of Sentinel-2 imagery has been constrained by the absence of ground-truth references for its \SI{20}{m} spectral bands,  forcing evaluations to depend on synthetic degradations or small cross-sensor proxies. We introduce SEN2NEON, a benchmark dataset that aligns high-resolution hyperspectral aerial imagery with Sentinel-2 reflectances, providing, for the first time, reliable reference data for these bands. This enables direct computation of standard SR metrics, as well as more sophisticated methods like the OpenSR-Test toolbox, to assess the multispectral integrity of SR models. Using SEN2NEON, we quantitatively compare two paradigm families: (i) \emph{frequency-transfer} models (SEN2SR variants) that inject \SI{10}{m}-derived spatial detail into the \SI{20}{m} targets following Wald’s scale-consistency principle, and (ii) \emph{direct prediction} via an 8$\times$ SRGAN. \textcolor{red}{The results show that the Wald protocol in combination with diffusion models produces sharper visual detail at the cost of hallucinations, whereas physically constrained models achieve superior spectral fidelity. SEN2NEON thus establishes, for the first time, a large-scale \SI{20}{m} Sentinel-2 SR benchmark under real reflectance conditions, enabling reproducible, evidence-based model selection for operational EO pipelines and providing a foundation for future multispectral SR research.}
\end{abstract}

\begin{IEEEkeywords}
super-resolution, sentinel-2, remote sensing, dataset
\end{IEEEkeywords}

\section{Introduction}
% General SR inrto

\iffalse
    Super-resolution (SR) has emerged as a powerful tool to enhance the effective spatial resolution of Earth Observation (EO) imagery~\cite{sr1,sr2}, allowing improved performance in applications such as crop field delineation~\cite{crop}, road extraction~\cite{road1} and building segmentation~\cite{building1,building2}. Sentinel-2, one of the most widely used EO data sources, provides multispectral imagery with heterogeneous spatial resolutions: \SI{10}{m} for RGB and NIR bands, \SI{20}{m} for red-edge and SWIR bands, and \SI{60}{m} for atmospheric correction bands~\cite{DRUSCH201225}.
    % General Sentence on most model focus on 10m bands, leaving the \SI{20}{m} bands untouched, breaking downstream approaches that rely n the standard rations
    Deep learning–based SR methods have recently shown great promise in recovering fine-scale detail~\cite{sr3}, in particular for the RGB bands \cite{wolters2023zooming,srgan1,Geltser2024}. Especially for the \SI{20}{m} bands, most published work relies on synthetic datasets created by artificially downsampling high-resolution imagery~\cite{Geltser2024}, or on limited cross-sensor collections such as SEN2VEN$\mu$S ~\cite{michel2022sen2venus} or NAIP-based datasets \cite{wolters2023zooming}. Synthetic training provides good spatial alignment, but fails to model real sensor noise and atmospheric effects, which often leads to learning the applied blurring kernel instead of real-world degradation \cite{opensrtest}. In contrast, cross-sensor datasets are usually small and geographically biased, limiting the statistical power of quantitative evaluation. Furthermore, these data sets suffer from spatial, temporal, and most importantly spectral misalignment due to differing acquisition dates, inaccuracies in georeferencing, and different spectral band windows.  
\fi
% General Intro + Val problems
Single-image super-resolution (SR) enhances the effective spatial resolution of Earth-observation (EO) imagery~\cite{sr1,sr2}, improving downstream tasks such as crop delineation~\cite{crop}, road extraction~\cite{road1}, and building segmentation~\cite{building1,building2}. Sentinel-2 provides 13 multispectral bands with heterogeneous ground sampling distances: \SI{10}{m} (RGB, NIR), \SI{20}{m} (red-edge, SWIR), and \SI{60}{m} (atmospheric)~\cite{DRUSCH201225}. Deep learning–based SR methods have shown strong performance in recovering fine-scale details~\cite{sr3}, especially for the \SI{10}{m} Sentinel-2 bands~\cite{wolters2023zooming,srgan1,Geltser2024}. However, work targeting the \SI{20}{m} bands remains limited. Most studies rely on synthetic datasets created by artificially downsampling high-resolution imagery~\cite{Geltser2024}, or on small cross-sensor collections such as SEN2VEN$\mu$S~\cite{michel2022sen2venus} and NAIP-based datasets~\cite{wolters2023zooming}. Synthetic training provides accurate spatial alignment, but fails to model real sensor noise and atmospheric effects, often leading networks to learn the degradation kernel rather than real-world image formation~\cite{opensrtest}. In contrast, cross-sensor datasets capture genuine sensor characteristics but are typically small, geographically biased, and prone to spatial, temporal, and spectral misalignment due to acquisition differences and georeferencing inaccuracies.

%While deep learning SR has advanced substantially for the \SI{10}{m} bands, progress on the \SI{20}{m} bands is hampered by the absence of high-resolution ground truth, forcing evaluations to rely on synthetic degradations or small cross-sensor datasets with spectral and temporal mismatches. This gap prevents rigorous, large-scale, quantitative benchmarking of Sentinel-2 \SI{20}{m} SR.

A major obstacle for the community has been the absence of high-resolution ground truth for the \SI{20}{m} Sentinel-2 bands, which prevents rigorous quantitative benchmarking and favors qualitative or proxy-based evaluation. To address this gap, we introduce SEN2NEON, a novel benchmark dataset created by harmonizing 1~m hyperspectral aerial imagery from NASA's AVIRIS-NG (Airborne Visible/Infrared Imaging Spectrometer - Next Generation) program with Sentinel-2 reflectances and rigorously aligning it to Level-2A product characteristics. This enables, for the first time, large-scale quantitative validation of Sentinel-2 SR models beyond the RGB-NIR bands.  

Beyond dataset construction, we demonstrate its utility through a comparative evaluation of representative SR paradigms. In this letter, we use SEN2NEON to compare two approaches: (i) direct prediction, where models reconstruct the \SI{20}{m} bands directly from the training dataset, and (ii) Wald-protocol-inspired frequency transfer, where high-resolution spatial frequencies extracted from the \SI{10}{m}-dervived SR product are used to guide the \SI{20}{m} band SR process. Their performance is evaluated using standard computer vision and remote-sensing metrics together with SR-specific indicators from the OpenSR-Test framework~\cite{opensrtest}.


\begin{figure*}[t]
    \centering
    \includegraphics[width=\textwidth]{sen2neon_viz.png}
    \caption{\textbf{SEN2NEON LR/HR image pair example.}
    Top row shows RGB of the original Sentinel-2 imagery at \SI{10}{m}, short-wave infrared at \SI{20}{m}, and color infrared at \SI{10}{m}.
    Bottom row shows the corresponding HR versions at \SI{2.5}{m} adapted from AVIRIS-NG.}
    \label{fig:sen2neon_lr_hr}
\end{figure*}

\section{SEN2NEON Dataset}

    \subsection{Data sources and coverage}
        SEN2NEON pairs high-resolution hyperspectral imagery from NASA's AVIRIS-NG (Airborne Visible/Infrared Imaging Spectrometer - Next Generation) program with temporally coincident Sentinel-2 Level-2A surface reflectance products. The AVIRIS-NG data, accessed through Google Earth Engine's NEON collections (projects/neon-prod-earthengine/assets/HSI\_REFL/001 and /002), provide 426 contiguous spectral bands spanning \SI{380}{nm}--\SI{2510}{nm} at native \SI{1}{m} spatial resolution. We filtered the collection to acquisitions after March 28, 2017, when both sensor stability and data availability significantly improved.
        
        For each AVIRIS-NG footprint, we identified temporally matched Sentinel-2 scenes within a $\pm$40-day window, prioritizing pairs with minimal temporal separation and high cloud confidence scores using Google's Cloud Score Plus product. The final dataset comprises 2,269 spatially aligned $2560\times2560$ pixel tiles at \SI{1}{m} resolution, spanning diverse land cover types across the continental United States: tree cover (1,186 tiles), grassland (581), shrubland (320), cropland (139), water (20), bare/sparse vegetation (19), and built-up areas (4). These tiles are derived from 127 unique AVIRIS-NG acquisition dates spanning 2017--2024, with temporal separations between AVIRIS-NG and Sentinel-2 acquisitions constrained to $\leq$40 days and cloud-free coverage exceeding 50\% in both sensors.
        
        
    \subsection{Spectral harmonization}
    
        Direct comparison between AVIRIS-NG hyperspectral data and Sentinel-2's multispectral bands requires rigorous spectral resampling. We convolved AVIRIS-NG's 426 bands (\SI{1}{nm} FWHM) with Sentinel-2's spectral response functions (SRFs) using linear interpolation to generate synthetic Sentinel-2 reflectances at \SI{1}{m} resolution. For each target Sentinel-2 band $b$, we compute
        
        \begin{equation}
        \rho_b^{\text{synth}} = \frac{\sum_{i} \text{SRF}_b(\lambda_i) \cdot \rho_i^{\text{AVIRIS}}}{\sum_{i} \text{SRF}_b(\lambda_i)}
        \end{equation}
        
        where $\rho_i^{\text{AVIRIS}}$ denotes AVIRIS-NG reflectance at wavelength $\lambda_i$, and the summation spans only wavelengths where $\text{SRF}_b(\lambda_i) > 0$. We applied Sentinel-2A and Sentinel-2B SRFs separately depending on the sensor identifier in each scene's metadata, accounting for known differences in the red-edge and SWIR sensor characteristics between the two satellites.
        
        This procedure yields 12 Sentinel-2-equivalent bands (B1--B9, B8A, B11--B12) at \SI{1}{m}. Notably, B10 (\SI{1375}{nm} cirrus band) is excluded from SEN2NEON. AVIRIS-NG wavelengths between \SI{1340}{nm}--\SI{1445}{nm} are set to invalid values (-100) due to atmospheric water vapor absorption, preventing spectral resampling for B10 centered at \SI{1375}{nm}. However, this omission does not compromise the dataset's utility: B10 is absent from Sentinel-2 L2A surface reflectance products, as it is used exclusively for atmospheric correction during Level-1C to Level-2A processing. Since SEN2NEON targets surface reflectance super-resolution benchmarking, the exclusion of B10 maintains consistency with operational L2A data products where this band is not available for analysis.
    
    
     \subsection{Reference generation and harmonization}
    
        From the native \SI{1}{m} synthetic Sentinel-2 imagery, we generated two reference resolutions via bilinear downsampling with antialiasing: \SI{2.5}{m} ($1024\times1024$ pixels) for $8\times$ SR evaluation and \SI{10}{m} ($256\times256$ pixels) for $2\times$ SR from \SI{20}{m} to \SI{10}{m}---the most operationally relevant task.
        
        To mitigate radiometric biases between AVIRIS-derived references and actual Sentinel-2 L2A products, we computed per-band linear correction coefficients via ordinary least squares regression:
        
        \begin{equation}
        \rho_{\text{S2}} = a \cdot \rho_{\text{AVIRIS-synth}} + b
        \end{equation}
        
        where $\rho_{\text{S2}}$ is Sentinel-2 L2A reflectance at \SI{10}{m} and $\rho_{\text{AVIRIS-synth}}$ is the downsampled synthetic reference. Fitted coefficients were applied to both resolution variants, with final values clipped to [0, 1], scaled by 10,000, and stored as 16-bit unsigned integers.
            
    



\iffalse
    \begin{itemize}
      \item Data sources: 1 m aerial multispectral imagery + Sentinel-2 L2A.
      \item Preprocessing: reflectance harmonization, PSF-based kernel degradation, geometric alignment.
      \item Coverage and statistics: area, temporal distribution, cloud filtering.
      \item Unique feature: enables true LR–HR quantitative testing for Sentinel-2 SR models.
    \end{itemize}
\fi


\begin{table*}[t]
\centering
\caption{Quantitative Performance Benchmarks for Sentinel-2 \SI{20}{m} Band Super-Resolution on the SEN2NEON Dataset}
\label{tab:validation_results}
\footnotesize
\setlength{\tabcolsep}{3pt}
\renewcommand{\arraystretch}{1.15}
\resizebox{1.03\textwidth}{!}{%
\begin{tabular}{lccccccccccc}
\toprule
\textbf{Model} &
\multicolumn{3}{c}{\textit{Standard Metrics}} &
\multicolumn{3}{c}{\textit{OpenSR Consistency}} &
\multicolumn{4}{c}{\textit{OpenSR Correctness}} \\
\cmidrule(lr){2-4} \cmidrule(lr){5-7} \cmidrule(lr){8-11}
 & PSNR ($\uparrow$) & SSIM ($\uparrow$) & SAM ($\downarrow$)
 & Refl. ($\downarrow$) & Spect. ($\downarrow$) & Spatial ($\downarrow$)
 & Synth. ($\uparrow$) & Omiss. ($\downarrow$) & Halluc. ($\downarrow$) & Improv. ($\uparrow$) \\
\midrule
\textbf{S2SR-LDSR-S2} & 31.98 $\pm$ 4.21 & 0.78 $\pm$ 0.12 & 0.05 $\pm$ 0.03
& 0.0087 $\pm$ 0.00 & 2.20 $\pm$ 1.33 & 0.019 $\pm$ 0.05
& 0.013 $\pm$ 0.01 & 0.40 $\pm$ 0.08 & 0.45 $\pm$ 0.08 & 0.16 $\pm$ 0.03 \\
\textbf{S2SR-Lite} & \textbf{32.56} $\pm$ 3.28 & 0.80 $\pm$ 0.10 & \textbf{0.04} $\pm$ 0.02
& \textbf{0.0054} $\pm$ 0.00 & 1.52 $\pm$ 0.91 & \textbf{0.004} $\pm$ 0.02
& 0.012 $\pm$ 0.01 & 0.47 $\pm$ 0.07 & \textbf{0.36} $\pm$ 0.08 & 0.18 $\pm$ 0.04 \\
\textbf{S2SR-Mamba} & 32.44 $\pm$ 4.18 & \textbf{0.81} $\pm$ 0.10 & 0.05 $\pm$ 0.03
& 0.0081 $\pm$ 0.00 & 2.20 $\pm$ 1.25 & 0.006 $\pm$ 0.04
& 0.013 $\pm$ 0.01 & 0.39 $\pm$ 0.08 & 0.42 $\pm$ 0.09 & \textbf{0.19} $\pm$ 0.04 \\
\textbf{SRGAN} & 31.40 $\pm$ 3.85 & 0.73 $\pm$ 0.13 & 0.05 $\pm$ 0.03
& 0.0067 $\pm$ 0.00 & \textbf{1.21} $\pm$ 0.79 & 0.020 $\pm$ 0.03
& \textbf{0.016} $\pm$ 0.01 & \textbf{0.35} $\pm$ 0.08 & 0.50 $\pm$ 0.09 & 0.16 $\pm$ 0.04 \\
\bottomrule
\end{tabular}
}
\vspace{2.5mm}
\begin{minipage}{0.95\textwidth}
\centering
\footnotesize
$(\uparrow)$ denotes a higher-is-better metric; $(\downarrow)$ lower-is-better. Best mean in \textbf{boldface}.
\end{minipage}
\end{table*}

\begin{figure*}[t]
    \centering
    \includegraphics[width=\textwidth]{grid_batch.png}
    \caption{Low resolution input, SR performance by model and the corresponding SEN2NEON reference for a random combination of the 20m S2 bands per row.}
    \label{fig:grid}
\end{figure*}


\section{Methodology}
\label{method}
We conduct our experiments on the SEN2NEON dataset, applying multiple super-resolution models and evaluating their outputs through a set of quantitative metrics to assess reconstruction quality and physical plausibility. We evaluate two SR families: Frequency-transfer (three SEN2SR variants \cite{sen2sr}: Mamba, Lite CNN, and a latent diffusion model \cite{donike}) that estimate high-resolution detail from the \SI{10}{m} bands and transfer it to the \SI{20}{m} targets, and a direct prediction SR-Generative Adversarial Network (SRGAN) that attempts to reconstruct the \SI{2.5}{m} detail directly from the six \SI{20}{m} bands. The 8$\times$ SRGAN is trained on $100k$ golabally sampled Sentinel-2 image patches, where the HR target corresponds to the native \SI{20}{m} bands and the LR input is synthetically degraded to \SI{160}{m} and subsequently upsampled back to \SI{20}{m} for network input. This setup emulates an 8$\times$ super-resolution task while maintaining alignment between LR and HR pairs. This workflow imitates similar synthetic dataset approaches. All models are evaluated on SEN2NEON using the same splits and preprocessing.

%\subsection{Evaluation Metrics}
%\label{metrics}

We evaluate the reconstructed imagery using both conventional image quality metrics and the domain-specific OpenSR-Test framework, providing a comprehensive assessment of pixel-level fidelity, spectral preservation, and physical correctness.

Standard metrics include PSNR, which measures radiometric accuracy through mean-squared error, SSIM for structural similarity~\cite{ssim}, and the Spectral Angle Mapper (SAM), which quantifies the angular difference between SR and HR spectra, providing an illumination-invariant measure of spectral preservation. Together, these metrics capture complementary aspects of radiometric accuracy, structural fidelity, and spectral consistency commonly used in computer vision and remote sensing. To complement these global measures, we adopt the OpenSR-Test framework~\cite{opensrtest}, designed for single-image SR in Earth observation. Its \textit{Consistency} metrics assess fidelity to the low-resolution input in terms of reflectance, spectral, and spatial agreement. \textit{Synthesis} quantifies the additional high-frequency information generated beyond a bilinear baseline, while \textit{Correctness} decomposes this detail into \textit{improvements}, \textit{omissions}, and \textit{hallucinations}, representing recovered, missing, and spurious structures, respectively. An ideal SR model maximizes improvements while minimizing omissions and hallucinations, providing a physically meaningful assessment of performance beyond pixel-wise similarity.

\iffalse
    Standard metrics include PSNR, which measures radiometric accuracy through mean-squared error, SSIM for structural similarity~\cite{ssim}, and the SAM, which quantifies the angular difference between SR and HR spectra, offering an illumination-invariant measure of spectral preservation. Together, these metrics capture complementary aspects of radiometric accuracy, structural fidelity, and spectral consistency that are widely used in both computer vision and remote sensing. To complement these global measures, we employ the OpenSR-Test framework~\cite{opensrtest}, specifically designed for single-image SR in Earth observation. OpenSR-test \textit{Consistency} metrics quantify whether the SR remains faithful to the original low-resolution input LR. We report reflectance consistency (mean absolute reflectance difference), spectral consistency (spectral angle distance between SR and LR spectra), and spatial consistency (spatial error of algorithmically derived key points), ensuring that radiometric and geometric properties are preserved during super-resolution. The \textit{Synthesis} metrics measure the amount of additional high-frequency information generated by the model beyond the bilinearly upsampled LR image. A higher synthesis score indicates that the model has introduced more detail, although not necessarily correctly. \textit{Correctness} then decomposes this added detail into three disjoint categories: improvements, omissions, and hallucinations. \textit{Improvements} correspond to high-frequency structures present in both SR and HR and therefore the desired recovery of information, \textit{Omissions} capture true details present in HR but missing from SR, and \textit{Hallucinations} quantify introduced frequencies that appear in SR but not in HR. An ideal SR model maximizes the improvements while minimizing both omissions and hallucinations. This framework thus provides a physically meaningful assessment of super-resolution performance, going beyond pixel-wise errors and enabling a deeper understanding of model behavior.
\fi

\section{Results}
%\subsection{Quantitative Benchmark}
The quantitative performance of the evaluated super-resolution models is summarized in \autoref{tab:validation_results}, following the validation metric pipeline described in Section~\ref{method}. Among the evaluated models, SEN2SR-Lite achieved the highest PSNR (32.56~dB) and lowest SAM, indicating superior radiometric accuracy and spectral preservation. The SEN2SR-Mamba variant reached the highest SSIM (0.81) and the strongest improvement score, confirming that its transformer-based architecture effectively enhances fine structural detail while maintaining fidelity. Both SEN2SR variants outperformed the diffusion-based S2SR-LD baseline, highlighting the advantage of lightweight and recurrent architectures over heavy generative models when trained under physical constraints. 

The SRGAN, which directly predicts the six \SI{20}{m} bands, produced the largest amount of added high-frequency content according to its synthesis score, but also exhibited the highest hallucination rate and the lowest SSIM, indicating that many of the introduced frequencies do not correspond to true structures. Reflectance and spectral consistency metrics further confirm this contrast: physically constrained SEN2SR models generate smoother, radiometrically coherent outputs, whereas SRGAN amplifies contrast and texture at the expense of spectral correctness and with limited correspondence to the actual spatial frequencies. Overall, these results demonstrate that frequency-transfer approaches preserve spectral integrity more effectively, and that spatial frequencies from the \SI{10}{m} bands are indeed correlated with those of the \SI{20}{m} bands, enabling their use to refine the SR products for these channels.


%\subsection{Qualitative Benchmark}
\autoref{fig:grid} provides a visual comparison of representative examples from the evaluated models on the SEN2NEON set. The SRGAN results appear noticeably sharper and more textured than the physically constrained models, consistent with its high synthesis and hallucination scores. However, these visual enhancements often correspond to artificial edges and exaggerated contrast not present in the reference imagery. In contrast, the SEN2SR variants produce spatially coherent reconstructions that preserve the underlying structures with fewer spectral artifacts, more closely resembling the high-resolution reference. The SEN2SR-Mamba outputs show particularly well-defined boundaries in heterogeneous areas, whereas the diffusion-based S2SR-LD tends to generate slightly smoother textures. Overall, the visual assessment supports the quantitative results: adversarial models emphasize perceptual sharpness, while frequency-transfer approaches maintain a more faithful spectral and structural representation of the scene.

\iffalse
    \begin{figure*}[t]
        \centering
        \includegraphics[width=\textwidth]{box_plot.png}
        \caption{
            \textbf{Model performance by Land Cover}.
        }
        \label{fig:sen2neon_lr_hr}
    \end{figure*}
\fi


%\subsection{Qualitative Results}
%BlaBlaBla looks great

\section{Discussion}
% General
While the presented experiments illustrate how SEN2NEON enables quantitative comparison of different super-resolution strategies, the dataset itself is the primary contribution of this work. By providing spatially and spectrally harmonized reference data for the \SI{20}{m} Sentinel-2 bands, SEN2NEON makes it possible for the first time to evaluate models under real reflectance conditions rather than synthetic degradations or cross-sensor proxies. This allows benchmarking of future and current architectures on a physically meaningful scale and supports the development of models optimized for operational Earth-observation pipelines. The experiments further indicate that the spatial frequencies learned during the \SI{2.5}{m} super-resolution process are strongly correlated with those present in the \SI{20}{m} Sentinel-2 bands. This suggests that fine-scale structures reconstructed at higher resolutions correspond to genuine high-frequency information that can be transferred across bands, rather than to artificial details. Consequently, frequency-transfer approaches that exploit this cross-resolution correlation can refine \SI{20}{m} products more effectively than models relying solely on direct prediction.

% limits
Despite these advantages, several limitations should be acknowledged. Temporal mismatches of up to several weeks between AVIRIS-NG and Sentinel-2 acquisitions may introduce residual differences in surface reflectance due to seasonal or illumination changes. The dataset currently only covers North American biomes and is therefore limited in its global representativeness. Furthermore, while spectral resampling and linear harmonization minimize sensor bias, remaining discrepancies in point-spread functions and calibration may slightly affect quantitative scores. These factors do not diminish the utility of the dataset but underline the importance of interpreting results as comparative rather than absolute. Future extensions should focus on including additional geographic regions and extending the harmonization process to different multi- and hyperspectral sensors and missions.

\section{Conclusion}
This letter introduced SEN2NEON, the first benchmark dataset enabling quantitative evaluation of super-resolution models for the \SI{20}{m} Sentinel-2 bands under real reflectance conditions. By harmonizing high-resolution AVIRIS-NG hyperspectral imagery with Sentinel-2 Level-2A characteristics, the dataset provides physically consistent references for both spectral and spatial assessment. The accompanying experiments demonstrate how SEN2NEON enables objective comparison of diverse SR strategies, confirming that high-resolution frequency information can support the reconstruction of the \SI{20}{m} bands. Beyond serving as a baseline for model evaluation, SEN2NEON establishes a foundation for reproducible benchmarking and evidence-based model selection in operational EO-SR pipelines.

\section*{Resources}
The dataset, models, and validation examples can be found at \url{github.com/ESAOpenSR/SEN2NEON}.





\bibliographystyle{IEEEtran}
\bibliography{references}

\end{document}
